\unitchapter{1}{5}{Mathematical Reasoning \& Aptitude}{p1-05}

\unitmeta{Expected questions: \textasciitilde5 (10 marks) · Target: 5/5 (pure
practice unit)}

\begin{sylbox}

\begin{itemize}
\tightlist
\item[$\square$]
  Types of reasoning
\item[$\square$]
  Number series, letter series, codes
\item[$\square$]
  Relationships (blood relations)
\item[$\square$]
  Classification
\item[$\square$]
  Mathematical aptitude: fraction, time \& distance, ratio, proportion, percentage, profit \& loss,
  interest \& discounting, averages
\end{itemize}

\end{sylbox}

\needspace{125pt}

\hypertarget{p1-05--1-number-series--pattern-checklist}{%
\section{Number Series --- Pattern Checklist}\label{p1-05--1-number-series--pattern-checklist}}

A number series hides a rule that generates each term from the ones before it. The quickest way to find it is
to test the simplest rules first --- constant differences, then constant ratios --- and only then look for
squares, alternating patterns or mixed operations. Writing the differences beneath the series takes a few
seconds and solves most questions.

\begin{figure}[!htbp]
\centering
\begin{tikzpicture}[node distance=5.5mm]
  \node[grey, minimum width=28mm] (A) {\textbf{Given series}};
  \node[diamondbox, fill=fsand, below=of A] (B) {differences\\constant?};
  \node[diamondbox, fill=fsand, below=of B] (C) {second differences\\constant?};
  \node[diamondbox, fill=fsand, below=of C] (D) {ratios\\constant?};
  \node[diamondbox, fill=fsand, below=of D] (E) {squares or\\cubes $\pm k$?};
  \node[diamondbox, fill=fsand, below=of E] (F) {alternate terms\\form a pattern?};
  \node[rose, text width=46mm, below=of F] (G) {mixed rule: $\times a + b$, primes, Fibonacci, factorials};
  \foreach \n/\t in {B/{arithmetic progression}, C/{quadratic pattern, e.g. $n(n+1)$}, D/{geometric progression},
      E/{$n^2 \pm k$ or $n^3 \pm k$}, F/{two interleaved series}} {
    \node[sage, text width=42mm, right=14mm of \n] (\n y) {\t};
    \draw[arr] (\n) -- node[lbl, above] {yes} (\n y);
  }
  \draw[arr] (A) -- (B);
  \foreach \a/\b in {B/C, C/D, D/E, E/F, F/G} \draw[arr] (\a) -- node[lbl, right] {no} (\b);
\end{tikzpicture}
\caption{A checklist for number series: test the simple rules first and move down only when they fail.}
\end{figure}


Examples:

\begin{itemize}
\tightlist
\item
  2, 6, 12, 20, 30, \textbf{42} → n(n+1)
\item
  3, 7, 15, 31, 63, \textbf{127} → ×2 + 1
\item
  1, 8, 27, 64, \textbf{125} → n³
\item
  2, 3, 5, 7, 11, \textbf{13} → primes
\item
  1, 1, 2, 3, 5, 8, \textbf{13} → Fibonacci
\item
  4, 9, 25, 49, 121, \textbf{169} → squares of primes
\item
  1, 2, 6, 24, 120, \textbf{720} → ×2, ×3, ×4 \ldots{} (factorials)
\end{itemize}

\needspace{155pt}

\hypertarget{p1-05--2-letter-series--coding}{%
\section{Letter Series \& Coding}\label{p1-05--2-letter-series--coding}}

Letter questions become number questions once each letter is replaced by its position in the alphabet.
Memorise the positions and their reverses so that no time is lost counting; the reference table below is
worth learning by heart.

Position table (memorise both directions):

\begin{center}
\small
\setlength{\tabcolsep}{3.2pt}
\begin{tabular}{l*{13}{c}}
\toprule
Letter & A & B & C & D & E & F & G & H & I & J & K & L & M \\
Position & 1 & 2 & 3 & 4 & 5 & 6 & 7 & 8 & 9 & 10 & 11 & 12 & 13 \\
Reverse & 26 & 25 & 24 & 23 & 22 & 21 & 20 & 19 & 18 & 17 & 16 & 15 & 14 \\
\midrule
Letter & N & O & P & Q & R & S & T & U & V & W & X & Y & Z \\
Position & 14 & 15 & 16 & 17 & 18 & 19 & 20 & 21 & 22 & 23 & 24 & 25 & 26 \\
Reverse & 13 & 12 & 11 & 10 & 9 & 8 & 7 & 6 & 5 & 4 & 3 & 2 & 1 \\
\bottomrule
\end{tabular}
\end{center}


Trick: \textbf{EJOTY} = 5, 10, 15, 20, 25. Opposite pairs sum to 27 (A↔Z, M↔N).

\begin{itemize}
\tightlist
\item
  Coding: If CAT → DBU (+1 each), then DOG → \textbf{EPH}.
\item
  If CAT = 24 (3+1+20), then DOG = 4+15+7 = \textbf{26}.
\item
  Reverse coding: CAT → XZG (opposite letters).
\end{itemize}

\needspace{155pt}

\hypertarget{p1-05--3-blood-relations}{%
\section{Blood Relations}\label{p1-05--3-blood-relations}}

Relationship puzzles are confusing in words and trivial on paper. Draw a small family tree as you read each
clause, marking gender and generation, and the answer can simply be read off.

Use a \textbf{family tree diagram} with symbols:

\begin{figure}[!htbp]
\centering
\begin{tikzpicture}[m/.style={ink, rectangle, fill=fslate, minimum size=9mm, font=\small\bfseries},
                    f/.style={ink, circle, fill=frose, minimum size=9.5mm, font=\small\bfseries}]
  \node[m] (C) at (0,0) {C}; \node[f] (D) at (2.6,0) {D};
  \draw[fgink, line width=0.6pt, double, double distance=1.4pt] (C) -- (D);
  \coordinate (mid) at (1.3,0);
  \draw[thinline] (mid) -- (1.3,-1.0);
  \node[m] (A) at (0,-2) {A}; \node[f] (B) at (2.6,-2) {B};
  \draw[thinline] (A.north) |- (1.3,-1.0) -| (B.north);
  \node[lbl, anchor=west] at (3.3,0) {married couple (double line)};
  \node[lbl, anchor=west] at (3.3,-2) {siblings share one horizontal bar};
  % legend
  \node[m, minimum size=5mm, font=\scriptsize] at (7.9,0.2) {}; \node[lbl, anchor=west] at (8.25,0.2) {male};
  \node[f, minimum size=5.5mm, font=\scriptsize] at (7.9,-0.5) {}; \node[lbl, anchor=west] at (8.25,-0.5) {female};
\end{tikzpicture}
\caption{"A is the brother of B; B is the daughter of C; C is the husband of D." Drawn as a family tree, the
answer is immediate: D is A's mother. Use squares for males and circles for females.}
\end{figure}


Coded relations: "P + Q means P is father of Q; P × Q means P is sister of Q" → draw step-by-step.

\needspace{227pt}

\hypertarget{p1-05--4-classification-odd-one-out}{%
\section{Classification (Odd one out)}\label{p1-05--4-classification-odd-one-out}}

Look for the property shared by all options but one: being prime, a perfect square, a multiple of a number,
a vowel, a member of a category. When two different properties seem to work, prefer the more specific one.
Find the common property in 3 of 4: primes, squares, multiples, vowels, same-category items.
Example: 121, 144, 169, \textbf{190} → 190 is not a perfect square.

\needspace{191pt}

\hypertarget{p1-05--5-arithmetic-aptitude--formula-bank}{%
\section{Arithmetic Aptitude --- Formula Bank}\label{p1-05--5-arithmetic-aptitude--formula-bank}}

Arithmetic questions in Paper I are not difficult, but they reward speed. Each formula below replaces several
lines of working; understand where it comes from once, then use it directly in the examination.

\needspace{125pt}

\hypertarget{p1-05--percentage}{%
\subsection{Percentage}\label{p1-05--percentage}}

Percentages express every quantity as a fraction of 100, which makes comparison easy. Successive changes do
\emph{not} simply add: a 20\% rise followed by a 20\% fall leaves you 4\% poorer, because the fall is taken on a
larger base.

\begin{formulatable}
$x\%$ of $y$ equals $y\%$ of $x$ & x\% \text{ of } y = y\% \text{ of } x \\
Net effect of successive changes of $a\%$ and $b\%$ (negative for decreases) & a + b + \frac{ab}{100} \\
Cut in consumption needed to offset a price rise of $r\%$ & \frac{r}{100+r}\times 100\% \\
Single discount equivalent to $d_1\%$ followed by $d_2\%$ & d_1 + d_2 - \frac{d_1 d_2}{100} \\
\end{formulatable}


\needspace{125pt}

\hypertarget{p1-05--profit--loss}{%
\subsection{Profit \& Loss}\label{p1-05--profit--loss}}

Profit and loss percentages are always calculated on the \textbf{cost price} unless the question says otherwise;
discounts are always calculated on the \textbf{marked price}.

\begin{formulatable}
Profit percentage (on cost price) & \text{Profit\%} = \frac{SP-CP}{CP}\times 100 \\
Selling price for a profit of $P\%$ & SP = CP\times\frac{100+P}{100} \\
Selling price after a discount of $d\%$ on marked price $M$ & SP = M\times\frac{100-d}{100} \\
Two articles at the same $SP$, one at $x\%$ gain and one at $x\%$ loss & \text{always a loss of } \frac{x^2}{100}\% \\
False weight: 900 g sold as 1 kg & \text{gain} = \frac{100}{900}\times 100 = 11.11\% \\
\end{formulatable}


\needspace{125pt}

\hypertarget{p1-05--simple--compound-interest}{%
\subsection{Simple \& Compound Interest}\label{p1-05--simple--compound-interest}}

Simple interest is earned only on the original principal; compound interest is also earned on interest
already added. The difference between them grows with time, and the two-year difference formula is the most
frequently used shortcut.

\begin{formulatable}
Simple interest & SI = \frac{P\,R\,T}{100} \\
Amount at compound interest & A = P\Bigl(1+\frac{R}{100}\Bigr)^{T},\qquad CI = A - P \\
$CI - SI$ over two years & P\Bigl(\frac{R}{100}\Bigr)^{2} \\
$CI - SI$ over three years & P\Bigl(\frac{R}{100}\Bigr)^{2}\Bigl(3+\frac{R}{100}\Bigr) \\
Sum doubles in $T$ years at simple interest & R = \frac{100}{T} \\
\end{formulatable}


\needspace{125pt}

\hypertarget{p1-05--ratio-proportion-averages}{%
\subsection{Ratio, Proportion, Averages}\label{p1-05--ratio-proportion-averages}}

A ratio compares two quantities; a proportion states that two ratios are equal. Averages smooth a set of
values into one representative number --- but note that the average \emph{speed} over equal distances is a
harmonic mean, not the simple average of the speeds.

\begin{formulatable}
Proportion & a:b = c:d \iff ad = bc \\
Mean proportional of $a$ and $b$ & \sqrt{ab} \\
Third proportional to $a$ and $b$ & \frac{b^2}{a} \\
Average & \frac{\text{sum of values}}{\text{number of values}} \\
One of $n$ values replaced and the average changes by $d$ & \text{the value changed by } n\,d \\
Average speed over two equal distances at $x$ and $y$ & \frac{2xy}{x+y}\quad(\text{not } \tfrac{x+y}{2}) \\
\end{formulatable}


\needspace{125pt}

\hypertarget{p1-05--time-speed-distance}{%
\subsection{Time, Speed, Distance}\label{p1-05--time-speed-distance}}

Nearly every motion problem --- trains, boats, races --- reduces to the single relation distance = speed ×
time, applied carefully with consistent units and, where two bodies move, with their \emph{relative} speed.

\begin{figure}[!htbp]
\centering
\begin{adjustbox}{max width=\linewidth}
\begin{tikzpicture}
  \filldraw[fill=fsand, draw=fgink, line width=0.55pt] (0,0) -- (3.6,0) -- (1.8,2.8) -- cycle;
  \draw[thinline] (0.62,0.98) -- (2.98,0.98) (1.8,0) -- (1.8,0.98);
  \node[font=\Large] at (1.8,1.6) {$D$};
  \node[font=\Large] at (0.95,0.45) {$S$};
  \node[font=\Large] at (2.65,0.45) {$T$};
  \node[align=left, anchor=west, font=\small] at (4.2,1.4) {$D = S\times T$\\[2pt]$S = D / T$\\[2pt]$T = D / S$\\[2pt]\footnotesize cover the unknown; read the rest};
  \node[slate, text width=36mm, anchor=west] (k) at (9.6,2.0) {km/h $\to$ m/s: $\times\,\frac{5}{18}$};
  \node[slate, text width=36mm, anchor=west] (m) at (9.6,0.9) {m/s $\to$ km/h: $\times\,\frac{18}{5}$};
\end{tikzpicture}
\end{adjustbox}
\caption{The distance–speed–time triangle and the two unit conversions that appear in almost every problem.}
\end{figure}


\begin{itemize}
\tightlist
\item
  Relative speed: same direction \textbf{(a − b)}, opposite direction \textbf{(a + b)}.
\item
  Train crossing pole: time = L / S; crossing platform: (L + P) / S.
\item
  Boats: downstream = u + v, upstream = u − v; still water speed = (D + U)/2, stream = (D − U)/2.
\end{itemize}

\needspace{125pt}

\hypertarget{p1-05--time--work}{%
\subsection{Time \& Work}\label{p1-05--time--work}}

Think of work as a rate: if A finishes a job in \emph{a} days, A does 1/\emph{a} of it each day. Rates of people (or
pipes) working together simply add.

\begin{formulatable}
$A$ alone takes $a$ days, $B$ alone $b$ days; together & \frac{ab}{a+b}\ \text{days} \\
Pipes and cisterns & \text{add rates of filling pipes, subtract emptying ones} \\
Men, days and hours per unit of work & \frac{M_1 D_1 H_1}{W_1} = \frac{M_2 D_2 H_2}{W_2} \\
\end{formulatable}


\needspace{95pt}

\hypertarget{p1-05--worked-examples}{%
\section{Worked Examples}\label{p1-05--worked-examples}}

\begin{enumerate}
\def\labelenumi{\arabic{enumi}.}
\tightlist
\item
  A price rises 20\% then falls 20\%. Net? → 20 − 20 − 400/100 = \textbf{−4\%} (loss).
\item
  ₹5000 at 10\% CI for 2 yrs: CI − SI = 5000 × (0.1)² = \textbf{₹50}.
\item
  A car goes at 40 km/h and returns at 60 km/h. Average speed = 2×40×60/100 = \textbf{48 km/h}.
\item
  A can finish in 10 days, B in 15. Together = 150/25 = \textbf{6 days}.
\item
  Two articles each sold at ₹990, one at 10\% profit, another at 10\% loss. → \textbf{1\% loss}.
\end{enumerate}

\needspace{131pt}

\hypertarget{p1-05--6-deeper-dive--more-question-types-with-shortcuts}{%
\section{Deeper Dive --- More Question Types with Shortcuts}\label{p1-05--6-deeper-dive--more-question-types-with-shortcuts}}

\needspace{95pt}

\hypertarget{p1-05--61-types-of-reasoning}{%
\subsection{Types of Reasoning}\label{p1-05--61-types-of-reasoning}}

\begin{itemize}
\tightlist
\item
  \textbf{Deductive}: general → specific; conclusion certain. \textbf{Inductive}: specific → general; conclusion probable.
\item
  \textbf{Abductive}: inference to the best explanation (a doctor diagnosing from symptoms).
\item
  \textbf{Analogical}: reasoning from similarity of two cases.
\end{itemize}

\needspace{125pt}

\hypertarget{p1-05--62-mixtures--alligation}{%
\subsection{Mixtures \& Alligation}\label{p1-05--62-mixtures--alligation}}

Alligation is a quick way of finding the ratio in which two ingredients at different prices (or
concentrations) must be mixed to give a mixture of a desired mean value.

\begin{figure}[!htbp]
\centering
\begin{tikzpicture}
  \node[slate, minimum width=26mm] (c) at (0,2.4) {cheaper $c$};
  \node[slate, minimum width=26mm] (d) at (6,2.4) {dearer $d$};
  \node[sage, minimum width=26mm] (dm) at (0,0) {$d-m$};
  \node[sage, minimum width=26mm] (mc) at (6,0) {$m-c$};
  \draw[thinline] (c.south east) -- (mc.north west);
  \draw[thinline] (d.south west) -- (dm.north east);
  \node[sand, minimum width=24mm] (m) at (3,1.2) {mean $m$};
  \node[align=center, font=\small] at (3,-0.95) {cheaper : dearer $= (d-m) : (m-c)$};
\end{tikzpicture}
\caption{The rule of alligation: subtract crosswise to obtain the mixing ratio.}
\end{figure}


\textbf{Example}: Rice at ₹40/kg and ₹60/kg mixed to get ₹52/kg → ratio = (60 − 52) : (52 − 40) = 8 : 12 = \textbf{2 : 3}.

Repeated dilution: from a vessel of x litres, y litres replaced by water n times → remaining pure = x(1 − y/x)ⁿ.

\needspace{227pt}

\hypertarget{p1-05--63-ages}{%
\subsection{Ages}\label{p1-05--63-ages}}

Translate words into equations. "A is twice as old as B; 10 years ago A was three times as old as B":
A = 2B, A − 10 = 3(B − 10) → 2B − 10 = 3B − 30 → \textbf{B = 20, A = 40}.

\needspace{161pt}

\hypertarget{p1-05--64-partnership}{%
\subsection{Partnership}\label{p1-05--64-partnership}}

Profit shared in ratio of \textbf{capital × time}. A invests ₹5000 for 12 months, B ₹6000 for 8 months → 60000 : 48000 = \textbf{5 : 4}.

\needspace{95pt}

\hypertarget{p1-05--65-discounting-bankers--true-discount}{%
\subsection{Discounting (banker\textquotesingle s \& true discount)}\label{p1-05--65-discounting-bankers--true-discount}}

\begin{formulatable}
True discount on amount $A$ due in $T$ years at $R\%$ & TD = \frac{A\,R\,T}{100 + R\,T} \\
Present worth & PW = A - TD \\
Banker's discount (simple interest on the amount) & BD = \frac{A\,R\,T}{100} \\
Banker's gain & BG = BD - TD = \text{SI on } TD \\
\end{formulatable}


\textbf{Example}: A = ₹1100 due in 1 year at 10\%: TD = 1100 × 10/110 = ₹100; PW = ₹1000; BD = ₹110; BG = ₹10.

\needspace{130pt}

\hypertarget{p1-05--66-codingdecoding-variants}{%
\subsection{Coding--Decoding Variants}\label{p1-05--66-codingdecoding-variants}}

\begingroup\small

\begin{longtable}[]{@{}
  >{\raggedright\arraybackslash}p{(\columnwidth - 2\tabcolsep) * \real{0.1856}}
  >{\raggedright\arraybackslash}p{(\columnwidth - 2\tabcolsep) * \real{0.4750}}@{}}
\toprule\noalign{}
\begin{minipage}[b]{\linewidth}\raggedright
\textbf{Type}
\end{minipage} & \begin{minipage}[b]{\linewidth}\raggedright
\textbf{Example}
\end{minipage} \\
\midrule\noalign{}
\endhead
\bottomrule\noalign{}
\endlastfoot
Letter shift & COMPUTER → DPNQVUFS (+1) \\
Reverse alphabet & GOOD → TLLW (A↔Z) \\
Position sum & BAD = 2 + 1 + 4 = 7 \\
Word reversal & READ → DAER \\
Symbol substitution & If + means ×, − means ÷: 8 + 2 − 4 = 8 × 2 ÷ 4 = \textbf{4} \\
\end{longtable}

\endgroup

\needspace{95pt}

\hypertarget{p1-05--67-fractions--comparison}{%
\subsection{Fractions \& Comparison}\label{p1-05--67-fractions--comparison}}

\begin{itemize}
\tightlist
\item
  To compare fractions quickly, cross-multiply: 5/7 vs 7/10 → 50 vs 49 → \textbf{5/7 is larger}.
\item
  Recurring decimals: 0.333\ldots{} = 1/3; 0.2727\ldots{} = 27/99 = 3/11.
\end{itemize}

\needspace{95pt}

\hypertarget{p1-05--68-worked-mixed-set}{%
\subsection{Worked Mixed Set}\label{p1-05--68-worked-mixed-set}}

\begin{enumerate}
\def\labelenumi{\arabic{enumi}.}
\tightlist
\item
  Successive discounts 20\% and 10\% = 20 + 10 − 2 = \textbf{28\%}.
\item
  A sum becomes ₹1331 in 3 years at 10\% CI → principal = 1331/1.331 = \textbf{₹1000}.
\item
  15 men finish work in 20 days; how many days for 25 men? 15 × 20 / 25 = \textbf{12 days}.
\item
  Boat: 16 km downstream in 2 h, 8 km upstream in 2 h → speeds 8 and 4 → still water \textbf{6 km/h}, stream \textbf{2 km/h}.
\item
  Average of first 50 natural numbers = (50 + 1)/2 = \textbf{25.5}.
\end{enumerate}

\needspace{95pt}

\hypertarget{p1-05--previous-year-questions-pyq-pattern}{%
\section{Previous Year Questions (PYQ pattern)}\label{p1-05--previous-year-questions-pyq-pattern}}

\begin{enumerate}
\def\labelenumi{\arabic{enumi}.}
\tightlist
\item
  Next term: 2, 5, 10, 17, 26, ? → \textbf{37} (n² + 1)
\item
  Next term: 1, 4, 9, 16, 25, ? → \textbf{36}
\item
  Missing: 3, 12, 48, ?, 768 → \textbf{192} (×4)
\item
  Next: 6, 11, 21, 36, 56, ? → \textbf{81} (diffs 5,10,15,20,25)
\item
  Letter series: AZ, BY, CX, ? → \textbf{DW}
\item
  If TEACHER is coded as VGCEJGT (+2), then STUDENT → \textbf{UVWFGPV}
\item
  If PEN = 35 (16+5+14), INK = ? → 9+14+11 = \textbf{34}
\item
  Pointing to a photograph, a man says, "Her mother is the only daughter of my mother." The man is the girl\textquotesingle s: \textbf{Maternal uncle}
\item
  A is B\textquotesingle s sister, C is B\textquotesingle s mother, D is C\textquotesingle s father, E is D\textquotesingle s mother. A is E\textquotesingle s: \textbf{Great-granddaughter}
\item
  Odd one out: 27, 64, 125, 144 → \textbf{144} (not a cube)
\item
  A shopkeeper marks 25\% above CP and allows 10\% discount. Gain\% = 1.25 × 0.9 = 1.125 → \textbf{12.5\%}
\item
  A sum doubles in 8 years at SI. Rate = \textbf{12.5\%}
\item
  A train 150 m long passes a pole in 15 s. Speed = 10 m/s = \textbf{36 km/h}
\item
  Average of 5 numbers is 20; one number removed, average becomes 18. Removed number = 100 − 72 = \textbf{28}
\item
  Ratio of A:B = 2:3, B:C = 4:5. A:B:C = \textbf{8:12:15}
\item
  If the population increases 10\% annually, from 10000 after 2 years: \textbf{12100}
\item
  Salary increased by 25\%; by what \% must it be reduced to restore? 25/125 × 100 = \textbf{20\%}
\item
  Two trains 120 m and 180 m at 50 \& 40 km/h opposite directions. Time to cross = 300 / (90×5/18) = 300/25 = \textbf{12 s}
\end{enumerate}

\textbf{More practice questions}

\begin{enumerate}
\def\labelenumi{\arabic{enumi}.}
\setcounter{enumi}{18}
\tightlist
\item
  Milk at ₹30/L mixed with water (free) to sell at ₹30 with 20\% profit. Ratio milk : water = \textbf{5 : 1} (cost price of mixture 25 → (25 − 0) : (30 − 25))
\item
  Next term: 1, 3, 7, 15, 31, ? → \textbf{63}
\item
  Next term: 3, 5, 9, 17, 33, ? → \textbf{65} (×2 − 1)
\item
  Find the odd one: 3, 5, 11, 14, 17 → \textbf{14} (not prime)
\item
  A and B invest ₹20,000 and ₹30,000 for 6 and 4 months; profit ratio = 120000 : 120000 = \textbf{1 : 1}
\item
  Ten years ago, father was 4 times his son\textquotesingle s age; now he is twice. Son\textquotesingle s present age: S − 10 = (2S − 10)/4 → \textbf{15 years}
\item
  If \textquotesingle SCHOOL\textquotesingle{} is coded as \textquotesingle RBGNNK\textquotesingle{} (−1), then \textquotesingle COLLEGE\textquotesingle{} is: \textbf{BNKKDFD}
\item
  Pointing to a lady, Rahul says, "She is the daughter of my grandfather\textquotesingle s only son." The lady is Rahul\textquotesingle s: \textbf{sister}
\item
  40\% of a number is 60. 75\% of it = \textbf{112.5}
\item
  Ratio of present worth to true discount for ₹1210 due in 2 years at 10\% SI: PW = 1210/1.2 = 1008.33; TD = 201.67 → \textbf{5 : 1}
\end{enumerate}

\begin{revbox}

\begin{itemize}
\tightlist
\item
  Net \% change = a + b + ab/100 · Same SP, ±x\% ⇒ loss x²/100 \%
\item
  CI−SI (2 yr) = P(R/100)² · Avg speed = 2xy/(x+y)
\item
  EJOTY = 5,10,15,20,25 · Opposite letters sum to 27
\end{itemize}

\end{revbox}
